% 5-page (excluding references) condensed draft for ECCV Workshop style.
% Based on eccv2022submission_rewritten.tex (user-provided).
\documentclass[runningheads]{llncs}

\usepackage{graphicx}
\usepackage{amsmath,amssymb}
\usepackage{bm}
\usepackage{url}

\title{Terrain Sensing with Smartphone Structured Light:\\
Topology-Constrained 2D-DTW for Grid Pattern Matching}

\author{Nobuaki Tanaka}
\institute{Meiji University\\ \email{ev250024@meiji.ac.jp}}

\begin{document}
\maketitle

\begin{abstract}
Small rovers often struggle on visually ``flat'' but subtly uneven terrain.
We study a low-cost structured-light approach using a single smartphone:
the screen projects a high-contrast grid and the front camera captures its deformation.
A key challenge is robustly matching the observed grid intersections to the ideal lattice under
perspective distortion and missing detections.
We propose a \emph{topology-constrained 2D dynamic time warping (2D-DTW)} scheme that enforces
global grid consistency while keeping computation lightweight.
Given the matched correspondences, we triangulate local 3D ground profiles from a single frame.
Experiments with an Android prototype show that the grid-based method is more stable than
feature-based baselines on low-texture surfaces.
\keywords{structured light \and grid matching \and dynamic time warping \and mobile sensing \and terrain reconstruction}
\end{abstract}

\section{Introduction}
Low-cost mobile robots and CanSat-style rovers frequently operate on uneven terrain where small bumps
or local tilts can destabilize locomotion.
Depth sensors (e.g., LiDAR) can be bulky or expensive; in contrast, structured-light triangulation
projects a known pattern and infers 3D geometry from its deformation.
Modern smartphones already provide an integrated display and camera, motivating a phone-only
structured-light sensor for ground profiling.

The main technical bottleneck is \emph{grid correspondence}:
the projected lattice is observed under perspective distortion, partial occlusion, and noise.
Naively applying 1D-DTW independently to rows/columns breaks the global lattice topology, while
general 2D-DTW for images is computationally heavy.
We therefore introduce a topology-constrained 2D-DTW tailored to rectangular grids:
it performs column-wise alignment under a global monotonicity constraint and can be paired with a
row-wise pass for full lattice matching.

\textbf{Contributions.}
(i) We describe a smartphone-only structured-light pipeline for extracting grid intersections and reconstructing
local terrain unevenness.
(ii) We propose a topology-constrained 2D-DTW formulation that yields globally consistent grid matching
with a simple dynamic program.

\section{Related work (brief)}
Two-dimensional warping via dynamic programming has been studied for image/matrix alignment, including
monotonic 2D warping of Uchida and Sakoe.
DTW variants for rotation/affine invariance have been explored in handwriting and symbol recognition.
Our setting differs in that the target is a \emph{known rectangular grid} and we exploit that topology to
restrict the warping and reduce cost.

\section{Method}
\subsection{System setup and grid extraction}
A smartphone is rigidly mounted facing downward.
The display renders a high-contrast rectangular grid; the front camera captures the pattern on the ground.
Each frame is processed by a lightweight pipeline: band-limited filtering, LoG enhancement, skeletonization,
and intersection detection, yielding a set of observed grid nodes in image coordinates.

\subsection{Topology-constrained 2D-DTW for grid matching}
Let the ideal display grid have $q$ vertical columns and $p$ rows of intersections.
From each frame, we sort detected intersections into \emph{column candidates} (roughly vertical groups)
and represent each observed column as a 1D sequence of row positions/feature values.
Denote by $A_i\in\mathbb{R}^{p}$ the $i$-th ideal column profile and by $B_j\in\mathbb{R}^{r}$ the $j$-th observed
column profile ($r$ may vary due to missing detections).

\paragraph{Step 1: column-to-column distance landscape.}
We compute a DTW distance for each pair of columns
\[
D_{i,j} := \mathrm{DTW}(A_i,B_j),\qquad i=1,\dots,q,\ \ j=1,\dots,s,
\]
forming a matrix $D\in\mathbb{R}^{q\times s}$.

\paragraph{Step 2: ``river path'' enforcing global consistency.}
We find a single monotone path through $D$ that follows a low-cost valley.
Formally, we minimize the cumulative cost along a DTW-style warping path on the $(i,j)$ grid.
This yields a globally consistent mapping $a(i)$ from ideal columns to observed columns.
We then repeat the same procedure on rows to obtain consistent row indices, and finally assign each detected
intersection a unique lattice index $(\hat{k},\hat{i})$.

\paragraph{Why this is ``2D'' while staying light.}
The 2D aspect is the global optimization over the \emph{matrix} $D$ (not independent 1D alignments),
which suppresses abrupt column jumps and preserves lattice topology.
The computation is still based on 1D DTW primitives and a 2D dynamic program over $D$.

\subsection{Single-frame triangulation}
\begin{figure}[t]
\centering
\includegraphics[width=0.92\linewidth]{smartphone_triangulatin.ng}
\caption{\textbf{Triangulation geometry in our smartphone structured-light setup.}
A display point on the screen (virtual light source) and the camera ray associated with a detected grid intersection define two rays.
Their intersection with the ground plane (via least-squares) yields the 3D point.
The baseline between screen and camera and the phone-to-ground distance govern depth sensitivity.}
\label{fig:triangulation_geometry}
\end{figure}

With camera calibration $P=K[R\mid \bm t]$, each matched pair consists of
(i) a display grid node $\bm S(u_d,v_d)$ (treated as a virtual light source on the screen plane),
and (ii) an observed image pixel $\bm u$ for the corresponding ground intersection.
Assuming a locally planar ground under the phone, we triangulate the ground point as the least-squares
intersection of (a) the camera back-projection ray from $\bm u$ and (b) the ray emitted from $\bm S(u_d,v_d)$
toward the ground plane.
The reconstructed set of intersections forms a dense local height map.
Figure~\ref{fig:triangulation_geometry} illustrates the geometric configuration used for reconstruction.

\section{Experiments}
\subsection{Setup}
We implemented the full pipeline as an Android application.
The phone was mounted at a fixed height ($h=3.0$\,cm) above the floor, and the viewing angle was restricted
to about $\pm 10^\circ$ using a narrow field-of-view film.
We compared: (i) the proposed structured-light grid + topology-constrained 2D-DTW matching + triangulation,
(ii) a feature-based baseline using ORB+RANSAC with conventional triangulation/bundle adjustment,
and (optionally) (iii) grid + nearest-neighbor matching without DTW regularization.
We tested multiple floor textures (high/medium/low) and placed blocks of known height for evaluation.


\begin{figure}[t]
\centering
\begin{minipage}[t]{0.48\linewidth}
\centering
\includegraphics[width=\linewidth]{20260228_20}
\end{minipage}\hfill
\begin{minipage}[t]{0.48\linewidth}
\centering
\includegraphics[width=\linewidth]{lower_half.ng}
\end{minipage}
\caption{\textbf{Prototype and raw input.}
\emph{Left:} our rover prototype with a downward-facing smartphone projecting a grid on the display.
\emph{Right:} a raw front-camera frame when the grid is projected onto an aluminum-foil surface; only a partial region of the pattern is visible, which currently leads to unstable correspondence estimation.}
\label{fig:prototype_and_raw}
\end{figure}

\paragraph{Practical note and current limitation.}
In early trials, the projected pattern sometimes appears only partially in the front camera (Fig.~\ref{fig:prototype_and_raw}, right), e.g., due to viewing angle, glare, or insufficient field of view.
Our current implementation assumes that a sufficient portion of the lattice is visible; when too few intersections are detected, the matching and triangulation become unreliable.
Improving robustness to partial observations (e.g., stronger priors on grid geometry, temporal integration, or automatic exposure control) is left for future work.

\paragraph{Why not use IR/FaceID depth?}
If a device provides an IR depth module (e.g., a structured-light or ToF sensor used for face authentication), obtaining depth can be much easier.
However, not all phones include such IR hardware, and additional depth pipelines may be computationally heavy for small robots.
Our approach targets the hardware-constrained setting: \emph{no IR camera}, \emph{commodity RGB camera}, and lightweight computation.
Moreover, the proposed topology-constrained 2D-DTW is a generic alignment tool and may be applicable beyond this specific grid pattern.

\paragraph{Ongoing experiments.}
We plan to continuously publish experiment videos, logs, and updates at: \url{https://www.youtube.com/@horoujhoma}.

\subsection{Observations}
On high-texture scenes, the feature baseline is competitive.
On medium/low-texture surfaces, the feature baseline often fails to produce enough stable correspondences,
while the structured-light grid keeps many valid matches and yields a more stable profile.
These results suggest that explicit structured-light correspondences are beneficial when natural features are scarce.

\section{Conclusion}
We presented a smartphone-only structured-light terrain sensing pipeline and a topology-constrained 2D-DTW
formulation for robust grid pattern matching.
The key idea is to enforce global lattice consistency by optimizing a monotone ``river path'' over a DTW distance
landscape, while keeping computation based on 1D DTW primitives.
Future work includes more extensive quantitative evaluation and integration with rover navigation.

\nocite{*}
\bibliographystyle{splncs04}
\bibliography{refs}
\end{document}